\documentclass[11pt,a4paper]{article}
\usepackage[margin=2.5cm]{geometry}
\usepackage{amsmath,amssymb}
\usepackage{booktabs}
\usepackage{graphicx}
\usepackage{caption}
\usepackage{microtype}
\usepackage[colorlinks=true,linkcolor=black,citecolor=black,urlcolor=black]{hyperref}

% Figures are made by make_figures.py into figures/. If one is missing, a box is shown instead.
\newcommand{\safeinclude}[2]{\IfFileExists{#1}{\includegraphics[width=#2]{#1}}%
  {\fbox{\parbox{#2}{\centering\vspace{2em}Missing figure: \texttt{\detokenize{#1}}\vspace{2em}}}}}

\newcommand{\ones}{\mathbf{1}}
\newcommand{\Sig}{\Sigma}

\title{Non-Normal returns and estimation error in mean--variance portfolios}
\author{Vash Ansari}
\date{\today}

\begin{document}
\maketitle

\begin{abstract}
\noindent
Daily stock returns have heavy tails, clustered volatility and correlations that rise in crises, yet
mean--variance optimisation assumes that means and covariances are known. Using 25 large US stocks and
the S\&P~500 from 2000 to 2025, I document these departures from Normality, derive the Markowitz
solution by Lagrange multipliers and verify it numerically, measure estimation error with random matrix
theory and a simulation, and compare 1/N, minimum-variance and shrinkage portfolios in a walk-forward
backtest with transaction costs. No optimised strategy reliably beat 1/N out of sample. Ledoit--Wolf
shrinkage and a long-only constraint both improved on the sample covariance, but the differences are
within the sampling noise of the Sharpe ratios.
\end{abstract}

\section{Question and data}
The question is how non-Normal real returns are, and whether mean--variance optimisation survives
estimation error. The data are daily adjusted closing prices (adjusted for splits and dividends)
downloaded from Yahoo Finance with \texttt{yfinance} on 11 October 2026 for 25 US large-cap
stocks and the S\&P~500 index: 6{,}538 prices from 3 January 2000 to 30 December 2025, hence 6{,}537
daily returns. Log returns are used for statistics and simple returns for portfolio arithmetic.
Means are annualised by 252 and volatilities by $\sqrt{252}$.

The data were checked before use. All 26 series are complete from the first date, no series has a run of
five or more identical prices, and there is one calendar gap longer than five days (17 September 2001,
after the post-attack market closure). About 64 daily moves exceed 20\% in absolute log return. They
cluster at known events, such as Apple on 29 September 2000 ($-0.73$), the banks in 2008--09 and several
stocks in March 2020, and none resembles an unadjusted split, so nothing was removed.

\section{Returns are not Normal}
\paragraph{Heavy tails.}
Excess kurtosis (zero for a Normal) is 114.6 for Apple and 89.0 for Procter \& Gamble, driven by single
days in 2000, and between 19 and 28 for BAC, WFC, HD, UNH and CVX; the S\&P~500 has 10.6. The
Jarque--Bera test rejects Normality for all 26 series. A Normal predicts 0.41 days beyond four standard
deviations per series in the sample; the data have between 27 and 64, about 65 to 155 times as many.
A Student-$t$ fitted by maximum likelihood has between 2.2 and 3.6 degrees of freedom for every series
(2.6 for the S\&P~500), and and Figure~\ref{fig:qq} shows it follows the bulk and moderate tails of the index far better than
a Normal, although the fit is imperfect in the far tails, where the few most extreme observations
are smaller than the fitted $t$ predicts.
With four or fewer degrees of freedom the fourth moment is infinite, so sample kurtosis is not a stable
quantity.

\begin{figure}[ht]
\centering
\safeinclude{figures/fig1_qq.png}{0.95\textwidth}
\caption{QQ plots of S\&P~500 daily log returns against a Normal (left) and a fitted Student-$t$ (right).}
\label{fig:qq}
\end{figure}

\paragraph{Volatility clustering.}
Returns themselves are almost uncorrelated: the average lag-1 autocorrelation is $-0.049$, which is
statistically detectable in 25 of 26 series with 6{,}537 observations but explains about 0.25\% of the
variance of the next return. Squared returns are very different: the average lag-1 autocorrelation is
0.193 and the Ljung--Box test (10 lags) rejects independence for all 26 series. Large moves follow large
moves; Figure~\ref{fig:clust} shows the slowly decaying autocorrelation of $|r|$ and the volatility
regimes of 2000--02, 2008--09 and 2020.

\begin{figure}[ht]
\centering
\safeinclude{figures/fig2_clustering.png}{\textwidth}
\caption{S\&P~500 rolling 21-day volatility (left) and autocorrelation of absolute returns (right).}
\label{fig:clust}
\end{figure}

\paragraph{Correlations in crises.}
The average pairwise correlation of the 25 stocks is 0.26 in 2006--07 and 0.30 in 2019, but 0.63 from
September 2008 to March 2009 and 0.72 from February to April 2020, while average volatility rises from
about 20\% to 72--80\%. Part of this increase may come from selecting high-volatility windows rather than
from a change in the underlying relationship \cite{forbes2002}, and the windows are short, so the figures
are indicative.

\section{Mean--variance optimisation}
Let $\mu$ be the vector of expected returns, $\Sig$ the covariance matrix and $w$ the weights. The
minimum-variance portfolio with expected return $m$ solves
\[
\min_w \tfrac12 w^\top \Sig w \quad\text{subject to}\quad \ones^\top w = 1,\ \ \mu^\top w = m .
\]
The Lagrangian is $L = \tfrac12 w^\top\Sig w - \lambda(\ones^\top w - 1) - \gamma(\mu^\top w - m)$, and
$\partial L/\partial w = \Sig w - \lambda\ones - \gamma\mu = 0$ gives
$w = \Sig^{-1}(\lambda\ones + \gamma\mu)$. With
$A = \ones^\top\Sig^{-1}\ones$, $B = \ones^\top\Sig^{-1}\mu$, $C = \mu^\top\Sig^{-1}\mu$ and
$D = AC - B^2$, substituting into the two constraints gives $\lambda A + \gamma B = 1$ and
$\lambda B + \gamma C = m$, so
\[
\lambda = \frac{C - Bm}{D},\qquad \gamma = \frac{Am - B}{D}.
\]
Without the return constraint ($\gamma = 0$, $\lambda = 1/A$) the global minimum-variance portfolio is
$w = \Sig^{-1}\ones/(\ones^\top\Sig^{-1}\ones)$. The variance along the frontier is
$\sigma^2(m) = w^\top\Sig w = \lambda + \gamma m = (Am^2 - 2Bm + C)/D$, a parabola in $m$, so the
frontier is a hyperbola in $(\sigma, m)$ with its leftmost point at $m = B/A$, $\sigma^2 = 1/A$. The
maximum-Sharpe (tangency) portfolio satisfies $\Sig w \propto (\mu - r_f\ones)$, hence
$w \propto \Sig^{-1}(\mu - r_f\ones)$ rescaled to sum to one. This is valid only when
$\ones^\top\Sig^{-1}(\mu - r_f\ones) > 0$, and the weights explode as that sum approaches zero.

The closed form was implemented with linear solves (never explicit inverses) and checked against a
numerical solver (SLSQP): the largest difference in any weight on the 25 stocks was $5.1\times10^{-6}$.
On the full sample, the minimum-variance portfolio has an annualised volatility of 14.9\%, below every
stock, with a return of 12.2\%. The tangency portfolio (with $r_f = 0$) is extreme: 10 negative weights,
weights from $-18\%$ to $+33\%$, and gross exposure of 3.05 (about 203\% long and 103\% short). Its
weights follow the estimated means, whose standard error is about six percentage points for a stock with
30\% volatility over 26 years, comparable to the whole spread of means in the sample.

\section{Estimation error}
\paragraph{Random matrix theory.}
For $n$ observations of $p$ independent variables with $q = p/n$, the eigenvalues of the sample
correlation matrix fall between $(1-\sqrt q)^2$ and $(1+\sqrt q)^2$ (Marchenko--Pastur
\cite{marchenko1967}). In the last 250 days ($q = 0.10$, upper edge 1.73) only two eigenvalues of the
stock correlation matrix lie above the edge (7.94 and 3.41); in the full sample (upper edge 1.13) three
do (9.31, 1.86 and 1.57). Everything inside the bulk is indistinguishable from noise
(Figure~\ref{fig:mp}).

\begin{figure}[ht]
\centering
\safeinclude{figures/fig3_mp.png}{0.7\textwidth}
\caption{Eigenvalues of the sample correlation matrix (full sample) against the Marchenko--Pastur density
for pure noise. The largest eigenvalue, the market factor, lies far to the right.}
\label{fig:mp}
\end{figure}

\paragraph{Shrinkage.}
The Ledoit--Wolf estimator \cite{ledoit2004} shrinks the sample covariance $S$ towards a scaled identity,
$\Sigma^\ast = \delta\, m I + (1-\delta) S$ with $m = \operatorname{tr}(S)/p$ and a data-driven
$\delta\in[0,1]$. My implementation matches scikit-learn exactly. The intensity is $\delta = 0.139$ with
250 days and 0.006 with the full sample, and the condition number falls from 109 to 34 in the short
window.

\paragraph{Simulation with a known truth.}
I drew data from a one-factor covariance matrix with 25 assets (500 repetitions per $n$) and, for each
estimate, built minimum-variance weights and measured their true variance relative to the best
achievable. Table~\ref{tab:sim} shows that estimation error costs 69\% extra variance at $n = 60$ with the
sample covariance, 34\% with Ledoit--Wolf, and falls towards zero as $n$ grows. The EWMA estimate
($\lambda = 0.94$, effective sample about 17 days) stays near 1.95 at every $n$ because its effective sample
is smaller than the number of assets and the simulated covariance is constant.

\begin{table}[ht]
\centering
\caption{True variance of minimum-variance portfolios relative to the optimum (1 is perfect).}
\label{tab:sim}
\begin{tabular}{rccc}
\toprule
$n$ & Sample & Ledoit--Wolf & EWMA \\
\midrule
60 & 1.689 & 1.335 & 1.984 \\
120 & 1.259 & 1.196 & 1.948 \\
250 & 1.109 & 1.098 & 1.941 \\
500 & 1.051 & 1.048 & 1.954 \\
1000 & 1.024 & 1.023 & 1.939 \\
\bottomrule
\end{tabular}
\end{table}

\section{Out-of-sample backtest}
\paragraph{Design.}
Each strategy is run walk-forward on the simple returns of the 25 stocks. Every 21 trading days the
weights are computed from the previous 504 days only (rows $t-504,\dots,t-1$ for a rebalance on day $t$)
and applied from day $t$; between rebalances they drift with returns. A transaction cost of 10 basis
points of turnover is deducted. The settings were fixed before looking at results. The strategies are 1/N;
minimum variance with the sample covariance; minimum variance with the Ledoit--Wolf covariance, with and
without a long-only constraint; and a tangency-type portfolio using Ledoit--Wolf covariance and means
shrunk halfway to their cross-sectional average ($k = 0.5$). The last falls back to minimum variance when
the tangency formula is invalid or gross exposure would exceed 3, which happened at 152 of 288
rebalances. Tests check that changing a future return cannot alter any earlier result.

\begin{table}[ht]
\centering
\caption{Backtest results, net of 10 bps costs (window 504 days, monthly rebalancing).}
\label{tab:main}
\small
\begin{tabular}{lrrrrrr}
\toprule
Strategy & Return & Vol. & Sharpe & Max DD & Turnover/yr & Gross Sharpe \\
\midrule
1/N                 & 14.6\% & 18.6\% & 0.823 & $-43.8\%$ & 0.58 & 0.826 \\
MinVar, sample      &  7.5\% & 14.5\% & 0.571 & $-40.2\%$ & 3.30 & 0.594 \\
MinVar, LW          &  8.0\% & 14.4\% & 0.609 & $-37.5\%$ & 2.83 & 0.629 \\
MinVar, LW long-only&  9.5\% & 14.5\% & 0.698 & $-35.3\%$ & 1.47 & 0.708 \\
Shrunk MV           & 11.0\% & 16.3\% & 0.725 & $-37.5\%$ & 8.43 & 0.777 \\
\bottomrule
\end{tabular}
\end{table}

\paragraph{Results.}
Equal weighting has the best Sharpe ratio (0.823) and ranks first in every sensitivity setting
(Table~\ref{tab:sens}), although it also has the deepest drawdown. The minimum-variance strategies lower
volatility from 18.6\% to about 14.5\% but give up more return than they save in risk. Ledoit--Wolf
beats the sample covariance in all four settings (by 0.03 to 0.06 of Sharpe) with lower turnover, and the
long-only constraint helps more still (0.698 against 0.609, half the turnover). The shrunk mean--variance
strategy has the highest turnover by far (8.4 a year), loses 0.05 of Sharpe to costs, and its rank moves
between second and fourth. In sub-periods, 1/N's lead comes largely from 2020--2025
(Table~\ref{tab:per}); in 2010--2019 all five strategies lie between 1.10 and 1.20.

\begin{figure}[ht]
\centering
\safeinclude{figures/fig4_backtest.png}{0.85\textwidth}
\caption{Walk-forward wealth and drawdown of the five strategies, net of costs.}
\label{fig:bt}
\end{figure}

\begin{table}[ht]
\centering
\caption{Net Sharpe ratio under alternative settings (window / rebalancing days / cost in bps).}
\label{tab:sens}
\small
\begin{tabular}{lccccc}
\toprule
Setting & 1/N & MinVar sample & MinVar LW & LW long-only & Shrunk MV \\
\midrule
252 / 21 / 10 & 0.805 & 0.566 & 0.625 & 0.698 & 0.621 \\
504 / 63 / 10 & 0.827 & 0.579 & 0.613 & 0.699 & 0.745 \\
504 / 21 / 0  & 0.826 & 0.594 & 0.629 & 0.708 & 0.777 \\
504 / 21 / 30 & 0.817 & 0.526 & 0.570 & 0.677 & 0.620 \\
\bottomrule
\end{tabular}
\end{table}

\begin{table}[ht]
\centering
\caption{Net Sharpe ratio by sub-period (main settings).}
\label{tab:per}
\small
\begin{tabular}{lccccc}
\toprule
Period & 1/N & MinVar sample & MinVar LW & LW long-only & Shrunk MV \\
\midrule
2002--2009 & 0.606 & 0.337 & 0.370 & 0.509 & 0.643 \\
2010--2019 & 1.188 & 1.122 & 1.157 & 1.204 & 1.100 \\
2020--2025 & 0.768 & 0.308 & 0.357 & 0.438 & 0.411 \\
\bottomrule
\end{tabular}
\end{table}

\paragraph{How much weight these differences carry.}
Over 23.9 years the approximate standard error of each Sharpe ratio is 0.22 to 0.24 (an i.i.d.
approximation). The lead of 1/N over the other strategies is 0.10 to 0.25, so at most about one standard
error. The strategies are correlated, which makes differences more precise than this, but I did not test
that, so the data show that nothing beat 1/N, not that 1/N is better.

\section{Conclusion and limitations}
Real returns have heavy tails, clustered volatility and crisis correlations, and these features make
sample means and covariances unreliable inputs. Shrinkage and constraints reduce the damage, in line with
the evidence that estimation error limits the gains from optimisation \cite{demiguel2009}, but a simple
equal-weight rule was not beaten here.

The limitations are the following. (i) \emph{Survivorship bias}: the stocks are today's large companies,
so failures are missing and results are flattered, possibly 1/N more than the others. (ii) One market, one
universe of 25 stocks and one 26-year period. (iii) The simulation assumes Normal data and a constant
one-factor covariance. (iv) Sharpe ratios have large standard errors. (v) Window, rebalancing, cost and
shrinkage constants were fixed in advance, and the sensitivity tables show how much rankings move.
(vi) Costs are proportional to turnover only: no market impact, borrowing cost for shorts or taxes, and
the risk-free rate is zero. (vii) The shrunk mean--variance strategy is partly minimum variance.

\begin{thebibliography}{9}
\bibitem{demiguel2009} V.~DeMiguel, L.~Garlappi and R.~Uppal. Optimal versus naive diversification: how
inefficient is the 1/N portfolio strategy? \emph{Review of Financial Studies}, 22(5):1915--1953, 2009.
\bibitem{forbes2002} K.~Forbes and R.~Rigobon. No contagion, only interdependence: measuring stock
market comovements. \emph{Journal of Finance}, 57(5):2223--2261, 2002.
\bibitem{ledoit2004} O.~Ledoit and M.~Wolf. A well-conditioned estimator for large-dimensional
covariance matrices. \emph{Journal of Multivariate Analysis}, 88(2):365--411, 2004.
\bibitem{marchenko1967} V.~A.~Marchenko and L.~A.~Pastur. Distribution of eigenvalues for some sets of
random matrices. \emph{Mathematics of the USSR-Sbornik}, 1(4):457--483, 1967.
\end{thebibliography}

\end{document}
